\section{Network Architecture and System Model}

\subsection{SD-VANET Architecture}

This research proposes a hierarchical computing architecture for task result caching in VANETs, consisting of four layers:

\begin{enumerate}
    \item \textbf{Vehicle Layer:} This layer includes all moving vehicles in the urban area, equipped with wireless communication devices to interact with RSUs and nearby vehicles. In this layer, there are $N$ different vehicles, identified by numbers $i = 1, 2, \ldots, N$, distributed across the network. Each vehicle generates computational tasks that may require offloading to fog infrastructure.
    
    \item \textbf{Fog Layer:} This layer comprises $M$ Roadside Units (RSUs), numbered $m = 1, 2, \ldots, M$, strategically positioned throughout the network. Each RSU is equipped with:
    \begin{itemize}
        \item \textbf{Computational Resources:} Processing capabilities for executing offloaded tasks from vehicles within coverage area.
        \item \textbf{Cache Storage:} Limited storage capacity $C_m$ for storing task results, enabling reuse of previously computed outputs.
        \item \textbf{Communication Interfaces:} Wireless links to vehicles, wired/wireless connections to other RSUs, and connection to the SDN controller.
    \end{itemize}
    
    \item \textbf{Cloud Layer:} This layer consists of powerful and stable resources located further from the network edge. Due to high communication latency, this layer is not suitable for latency-sensitive vehicular tasks but can serve as backup for delay-tolerant computations.
    
    \item \textbf{Control Layer:} Comprising a central SDN controller, this layer maintains a comprehensive view of network conditions, task distributions, cache states, and resource availability across all RSUs. The controller coordinates cache management decisions, including FL-based pre-placement and RL-based eviction policies.
\end{enumerate}

This architecture optimizes task processing and caching by leveraging distributed fog resources under centralized SDN control. Task $T_{ij}$, generated by vehicle $i$ with task identifier $j$, can be processed through one of the following mechanisms:

\begin{itemize}
    \item \textbf{Cache Hit:} If a matching task result exists in the RSU's cache, it is immediately returned without re-execution, dramatically reducing response time, energy consumption, and computational overhead.
    
    \item \textbf{Local Fog Execution:} If no cached result exists and the RSU has sufficient computational resources, the task is executed locally at the fog node, and the result is stored in cache for future reuse.
    
    \item \textbf{Remote Fog Execution:} If the local RSU lacks resources, the SDN controller redirects the task to another RSU with available capacity, balancing load across the network.
    
    \item \textbf{Cloud Offloading:} If fog resources are exhausted and task deadlines permit, the task is offloaded to the cloud layer for execution.
\end{itemize}

The SDN controller, with comprehensive network awareness, coordinates the FL-RL-TRC caching framework. The FL advisor analyzes task patterns across all RSUs and determines which results to pre-place in caches before requests arrive. The RL agent monitors real-time cache states and makes dynamic eviction decisions when storage capacity is reached. By integrating machine learning with SDN control, the system achieves intelligent cache management that adapts to dynamic vehicular patterns while maintaining global optimization objectives.

The network topology consists of multiple coverage regions, each served by an RSU with computational and caching capabilities. RSUs communicate with vehicles via wireless links and with each other through the SDN controller. Fig. \ref{fig:fl_rl_architecture} illustrates the FL-RL-TRC architecture, showing the interaction between vehicles, RSUs, the FL advisor, the RL agent, and the SDN controller.

\begin{figure}[t]
    \centering
    \includegraphics[width=\linewidth]{fl_rl_architecture.pdf}
    \caption{The FL-RL-TRC architecture for intelligent task result caching in SD-VANETs.}
    \label{fig:fl_rl_architecture}
\end{figure}

In this system, each vehicle first sends its task ($T_{ij}$) to the nearest RSU. The RSU checks its local cache for a matching result using a hash-based key derived from task intrinsics (task type, priority, data size). Upon cache hit, the result is immediately returned. Upon cache miss, the RSU forwards the task to the SDN controller, which schedules execution at an appropriate fog node based on resource availability and load balancing objectives. After execution, the result is stored in the cache and returned to the requesting vehicle. The FL advisor periodically analyzes task patterns and updates cache pre-placement decisions, while the RL agent continuously monitors cache states and makes eviction decisions when capacity limits are reached.

\subsection{Task Model}

Each task $T_{ij}$ generated by vehicle $i$ is characterized by the following attributes:

\begin{itemize}
    \item \textbf{Task Type} ($\tau_{ij}$): Category identifier from set $\{0, 1, 2, 3\}$ representing different computational requirements (e.g., route optimization, traffic analysis, sensor processing, hazard detection).
    
    \item \textbf{Priority} ($p_{ij}$): Urgency level from set $\{0, 1, 2, 3\}$ with higher values indicating greater importance for safety-critical applications.
    
    \item \textbf{Data Size} ($s_{ij}$): Input payload size in kilobytes, affecting both transmission time and cache storage requirements.
    
    \item \textbf{Deadline} ($d_{ij}$): Hard time constraint by which task must complete. Tasks missing deadlines are considered failed.
    
    \item \textbf{Computational Requirement} ($c_{ij}$): CPU cycles needed for task execution, measured in megacycles.
    
    \item \textbf{Origin RSU} ($o_{ij}$): The RSU where task was submitted.
    
    \item \textbf{Destination RSU} ($dest_{ij}$): The RSU where vehicle expects to receive result, accounting for mobility.
\end{itemize}

Tasks are generated stochastically by vehicles with probability $\lambda$ per timestep. The task generation follows realistic patterns where certain task types are more frequent in specific geographic regions and time periods. For cache key generation, we use a hash function based on task intrinsics:

\begin{equation}
    key_{ij} = \text{hash}(\tau_{ij}, p_{ij}, \lfloor s_{ij}/100 \rfloor)
\end{equation}

This approach enables cache reuse across different vehicles requesting identical computations, as the cache key depends only on task characteristics, not vehicle identity.

\subsection{Cache Model}

Each RSU $m$ maintains a cache with the following properties:

\begin{itemize}
    \item \textbf{Capacity} ($C_m$): Maximum number of task results that can be stored simultaneously.
    
    \item \textbf{Cache Entries}: Each entry stores $(key, result, metadata)$ where metadata includes:
    \begin{itemize}
        \item Last access time (for LRU tracking)
        \item Access frequency (for LFU tracking)
        \item Time-to-live (TTL) value
        \item Task characteristics (type, priority, size)
    \end{itemize}
    
    \item \textbf{Eviction Policies}: The RL agent can select from five actions:
    \begin{itemize}
        \item $a_0$: No operation (maintain current state)
        \item $a_1$: Evict least recently used entry
        \item $a_2$: Evict random entry
        \item $a_3$: Evict least frequently used entry
        \item $a_4$: Shrink TTL (halve TTL of all entries)
    \end{itemize}
\end{itemize}

Cache operations include:

\begin{equation}
    \text{Hit}(T_{ij}) = 
    \begin{cases}
        \text{True}, & \text{if } key_{ij} \in \text{Cache}_m \\
        \text{False}, & \text{otherwise}
    \end{cases}
\end{equation}

The cache hit rate is defined as:

\begin{equation}
    \text{HitRate} = \frac{\sum_{t=1}^{T} \text{Hits}_t}{\sum_{t=1}^{T} (\text{Hits}_t + \text{Misses}_t)}
\end{equation}

where $T$ is the total number of timesteps in the simulation.

\subsection{Delay Model}

To determine task execution delay within our architecture, we analyze three distinct scenarios:

\begin{itemize}
    \item \textbf{Cache Hit}: If the task result exists in cache, the delay consists only of communication time:
    \begin{equation}
        t_{ij}^{\text{cache}} = t_{\text{trans},ij} + t_{\text{recv},ij}
    \end{equation}
    where $t_{\text{trans},ij}$ is the transmission time from vehicle to RSU, and $t_{\text{recv},ij}$ is the return time from RSU to vehicle.
    
    \item \textbf{Local Fog Execution}: For cache miss with local execution, the total time is:
    \begin{equation}
        t_{ij}^{\text{fog}} = t_{\text{trans},ij} + \frac{c_{ij}}{r_m} + t_{\text{cache\_store}} + t_{\text{recv},ij}
    \end{equation}
    where $r_m$ is the computational capacity of RSU $m$, and $t_{\text{cache\_store}}$ is the time to store result in cache (typically negligible).
    
    \item \textbf{Remote Fog Execution}: For execution at remote RSU $m'$:
    \begin{equation}
        t_{ij}^{\text{remote}} = t_{\text{trans},ij} + t_{\text{forward},mm'} + \frac{c_{ij}}{r_{m'}} + t_{\text{return},m'i} + t_{\text{recv},ij}
    \end{equation}
    where $t_{\text{forward},mm'}$ is inter-RSU communication time and $t_{\text{return},m'i}$ is result return time.
\end{itemize}

The transmission time from vehicle $i$ to RSU $m$ is calculated as:

\begin{equation}
    t_{\text{trans},ij} = \frac{s_{ij}}{B \log_2(1 + \text{SNR}_{im})}
\end{equation}

where $B$ is the channel bandwidth and $\text{SNR}_{im}$ is the signal-to-noise ratio between vehicle $i$ and RSU $m$.

\subsection{Energy Consumption Model}

The energy consumption for task processing comprises two main components:

\begin{equation}
    E_{ij} = E_{\text{comm},ij} + E_{\text{comp},ij}
\end{equation}

The communication energy for transmitting task and receiving result is:

\begin{equation}
    E_{\text{comm},ij} = P_{\text{tx}} \cdot (t_{\text{trans},ij} + t_{\text{recv},ij})
\end{equation}

where $P_{\text{tx}}$ is the transmission power.

For cache hits, computational energy is zero since no execution occurs:

\begin{equation}
    E_{\text{comp},ij} = 
    \begin{cases}
        0, & \text{if cache hit} \\
        \kappa \cdot c_{ij}, & \text{if cache miss}
    \end{cases}
\end{equation}

where $\kappa$ is the energy coefficient per CPU cycle at the fog node.

By preventing redundant task execution through intelligent caching, FL-RL-TRC significantly reduces both communication and computational energy consumption. Cache hits eliminate the need for task execution, conserving fog node energy for other critical tasks and improving overall system efficiency.

\subsection{Performance Metrics}

The system performance is evaluated using the following metrics:

\begin{itemize}
    \item \textbf{Cache Hit Rate}: Percentage of tasks served from cache without re-execution.
    
    \item \textbf{Task Success Rate}: Percentage of tasks completed before their deadlines.
    
    \item \textbf{Average Completion Time}: Mean time from task submission to result delivery.
    
    \item \textbf{Energy Savings}: Reduction in computational energy compared to no-caching baseline.
    
    \item \textbf{System Throughput}: Number of tasks successfully completed per unit time.
\end{itemize}

The optimization objective is to maximize cache hit rate and task success rate while minimizing average completion time and energy consumption, subject to cache capacity constraints and hard deadline requirements.
